%==============================================================================
% Capítulo 10: Recomendaciones y líneas de trabajo futuro
%==============================================================================
\section{Recomendaciones y líneas de trabajo futuro}\label{sec:recomendaciones}

A partir de los resultados obtenidos, del análisis de errores presentado en la \autoref{subsec:calibracion} y de las limitaciones reconocidas en la \autoref{subsec:limitaciones-estudio}, se formulan las siguientes recomendaciones para la evolución del sistema y para futuras replicaciones del estudio.

Las recomendaciones que siguen no revisten idéntica prioridad. Se priorizan las líneas que atacan los defectos medidos y reconocidos entre las limitaciones de la \autoref{subsec:limitaciones-estudio} ---calibración del gate de confianza, clasificador supervisado con corpus propio y corrección de la medición de intervención humana---, por sobre las extensiones funcionales ---canales adicionales, paneles de monitoreo y base de conocimientos---, que amplían el alcance sin alterar la validez de lo medido.

\subsection{Incorporación de un clasificador supervisado con corpus propio}\label{subsec:clasificador-supervisado}

Se sugiere incorporar progresivamente un clasificador supervisado entrenado con el corpus etiquetado generado por el propio sistema durante su operación. La acumulación sostenida de decisiones validadas por los sectores responsables constituye un activo informacional valioso que, una vez alcanzado un volumen del orden de 5.000 casos (de carácter indicativo, no derivado del corpus actual), permitiría entrenar modelos de clasificación específicos del dominio organizacional. Las ventajas potenciales de esta migración incluyen: reducción del costo marginal por incidente clasificado al eliminar la dependencia del servicio externo de inferencia, disminución de la latencia al ejecutar la inferencia localmente, y ganancia de soberanía sobre los datos al evitar la transferencia internacional de fragmentos textuales. \textcite{Touvron2023_llama2} y trabajos subsiguientes sobre modelos de código abierto sugieren que los LLM desplegados localmente alcanzan rendimientos competitivos en tareas de clasificación cuando se invocan mediante prompts adecuadamente diseñados, lo que abre una vía concreta hacía un sistema completamente autoalojado.

\subsection{Ampliación de canales de entrada}\label{subsec:ampliacion-canales}

Se propone ampliar el conjunto de canales de entrada para incluir mensajería corporativa instantánea, particularmente Slack y Microsoft Teams en aquéllas organizaciones donde estas plataformas constituyen el canal de comunicación dominante. La inclusión de aplicaciones móviles propias para usuarios de campo constituye una extensión natural que atiende la demanda creciente de canales asincrónicos integrados al ecosistema laboral cotidiano. La arquitectura modular propuesta en la \autoref{sec:arquitectura} facilita esta ampliación: cada nuevo canal requiere únicamente un disparador adicional en N8N, sin modificaciones en las capas de procesamiento ni de persistencia.

\subsection{Panel de monitoreo en tiempo real}\label{subsec:panel-monitoreo}

El sistema incorpora un dashboard de analítica basado en ReCharts que permite visualizar tendencias históricas, distribución por sector y desglose por estado (véase \autoref{subsec:dashboard-analitica}). Como evolución de esta capacidad, se sugiere implementar una capa de monitoreo en tiempo real sustentada sobre tecnologías de observabilidad maduras, tales como Prometheus para la captura de métricas y Grafana para la visualización. Esta capa permitiría supervisar la latencia del sistema, la tasa de aciertos del clasificador, la distribución de carga entre sectores y el costo agregado de las invocaciones al modelo externo, habilitando alertas tempranas frente a degradaciones del servicio. La existencia de endpoints de salud (\texttt{GET /api/v1/health}) en cada componente del sistema proporciona los puntos de datos necesarios para instrumentar esta capa de observabilidad.

Más alla del monitoreo operativo, los datos acumulados de incidentes constituyen una fuente de inteligencia organizacional: su registro sistemático habilita flujos ETL (\emph{Extract, Transform, Load}) que permiten identificar patrones recurrentes de fallas, detectar áreas con mayor demanda de soporte y generar indicadores trimestrales que orienten decisiones estratégicas de inversión.

\subsection{Mecanismos de aprendizaje activo}\label{subsec:aprendizaje-activo}

Se propone incorporar mecanismos de aprendizaje activo (\emph{active learning}), en virtud de los cuales aquellos casos clasificados con baja confianza se prioricen automáticamente para revisión humana, y los resultados de esa revisión retroalimenten al sistema mediante un mecanismo de ajuste incremental. Este enfoque maximiza el rendimiento marginal de cada intervención humana ---que de todos modos ocurre en los casos derivados a revisión--- y acelera la mejora continua del clasificador con un costo operativo controlado.

\subsection{Estudios de replicación}\label{subsec:estudios-replicacion}

Se sugiere realizar estudios de replicación en organizaciones de distinto tamaño, sector y región geográfica, con el objeto de fortalecer la validez externa de los hallazgos y construir un cuerpo de evidencia empírica sobre el desempeño de soluciones similares en contextos heterogeneos. Particularmente relevantes serían las réplicas en organizaciones del sector público (con mayor proporción de incidentes vinculados a sistemas administrativos), en organizaciones del sector industrial (con mayor proporción de incidentes vinculados a infraestructura física) y en organizaciones de mayor tamaño donde el volumen diario de incidentes supere las 100 unidades.

\subsection{Evaluación de modelos de lenguaje de código abierto}\label{subsec:evaluacion-modelos-abiertos}

Se propone evaluar la integración de modelos de lenguaje de código abierto desplegados localmente ---tales como Llama~3, Mistral o Qwen--- como alternativa al modelo externo Gemini~3.6 Flash. Esta evaluación permitiría cuantificar el costo en exactitud frente a la ganancia en soberanía de datos, la eliminación de la dependencia de un proveedor único y la reducción de costos operativos a largo plazo. La realización de un benchmark comparativo sobre el mismo corpus de 200 casos constituiría una contribución adicional significativa al estado del arte.

\subsection{Base de conocimientos para resolución automática}\label{subsec:base-conocimientos}

El sistema desarrollado en este trabajo se concentra en la clasificación y derivación de incidentes hacía los sectores responsables, pero no aborda la resolución de aquellos casos que admiten soluciones estandarizadas y repetibles. Se recomienda, como línea de trabajo futuro, la incorporación de una base de conocimientos que permita al sistema no solo categorizar el incidente sino también proponer una resolución cuando el problema detectado coincida con patrones previamente documentados y validados.

La base de conocimientos operaría como un repositorio estructurado de soluciones verificadas. Ante un nuevo incidente, el sistema consultaría la base tras la clasificación inicial y, en caso de hallar una coincidencia de alta confianza, ofrecería al usuario los pasos necesarios para resolver el problema sin intervención del personal técnico. Este mecanismo de autoservicio reduciría la carga sobre los sectores de soporte al absorber los casos simples y repetitivos ---tales como reinicios de sesión, restablecimiento de contraseñas o configuraciones básicas de correo electrónico--- que representan una proporción significativa del volumen de incidentes en entornos de mesa de ayuda.

La construcción de esta base de conocimientos puede articularse con el mecanismo de aprendizaje activo descrito en la \autoref{subsec:aprendizaje-activo}: cada caso resuelto por un operador humano y validado como correcto constituye un candidato para ser incorporado a la base de conocimientos, previa curación editorial que garantice la claridad y la precisión de las instrucciones. La viabilidad técnica de este enfoque se ve reforzada por la creciente capacidad de los modelos de lenguaje para generar respuestas estructuradas a partir de documentación técnica, línea de investigación que conecta de manera directa con la evaluación de modelos de código abierto propuesta en la \autoref{subsec:evaluacion-modelos-abiertos}.

\subsection{Agente conversacional telefónico en tiempo real}\label{subsec:agente-conversacional}

Se propone, como línea de evolución de mayor alcance para el canal telefónico, reemplazar el esquema actual de captura de toma única ---una grabación de hasta 45 segundos que se transcribe de forma diferida--- por un agente que mantenga un diálogo en tiempo real con el usuario. Sobre Twilio Media Streams, que habilita el envío bidireccional del audio de la llamada hacia un servicio propio mediante WebSocket, el agente podría formular exclusivamente preguntas de aclaración pertinentes al problema reportado cuando la descripción inicial resultara insuficiente, consultar al usuario si desea agregar información tras detectar una pausa prolongada, comunicarle el número de incidente asignado una vez completado el registro, y cerrar la llamada de forma autónoma. Esta capacidad elevaría la calidad de la captura telefónica y reduciría la necesidad de contacto posterior, en línea con el tercer objetivo específico de ampliar la accesibilidad de los canales sin incrementar la complejidad percibida por el usuario final.

La adopción de esta arquitectura estuvo condicionada por el estado de madurez de las interfaces de audio bidireccional en tiempo real disponibles durante el desarrollo de este trabajo, que se encontraban en fase de vista previa (\emph{preview}) y con una superficie de programación sujeta a cambios, lo que representaba un riesgo de cronograma incompatible con los plazos del estudio. A ello se suma que la implementación exige infraestructura no contemplada en la arquitectura actual ---un punto de entrada WebSocket con soporte de \emph{upgrade} en el proxy inverso, la transcodificación del audio telefónico (mu-law, 8~kHz) hacia el formato de entrada del modelo (PCM, 16~kHz) y su retorno, la detección de actividad de voz y la gestión de turnos de habla--- junto con la redefinición del modelo de costos, hoy expresado por grabación, hacia un esquema por minuto de transmisión. Por estos motivos se recomienda abordar esta línea mediante una prueba de concepto acotada y conservar el mecanismo actual de grabación y transcripción como modo degradado de respaldo para garantizar la continuidad del canal.

\newpage
